% TOP_PROBLEM: {"id": 3402, "title": "P vs. PSPACE", "category_id": 15, "category": {"id": 15, "name": "computer_science", "display_name": "Computer Science", "description": "Computational complexity, algorithms, and theoretical CS.", "slug": "computer-science", "order_index": 15, "created_at": "2026-07-31T15:26:25.671Z"}, "status": "open"}
% FIRST_POSTED: null
% !TeX program = lualatex
% Compile twice: lualatex open_problems_500.tex
% All references and prose are embedded; no BibTeX database or external inputs.
\documentclass[11pt,a4paper]{article}
\usepackage[margin=24mm,headheight=14pt]{geometry}
\usepackage{fontspec}
\setmainfont{Latin Modern Roman}
\setsansfont{Latin Modern Sans}
\setmonofont{Latin Modern Mono}
\usepackage{amsmath,amssymb,mathtools}
\usepackage{unicode-math}
\setmathfont{Latin Modern Math}
\usepackage{microtype}
\usepackage{enumitem,xurl,needspace}
\usepackage[dvipsnames]{xcolor}
\usepackage{hyperref}
\hypersetup{unicode=true,colorlinks=true,linkcolor=MidnightBlue,urlcolor=MidnightBlue,
 pdftitle={Mathematical Research Notebook},pdfauthor={Source-annotated compilation},bookmarksnumbered=true}
\usepackage{bookmark}
\usepackage{tocloft}
\setlength{\cftsecnumwidth}{3em}
\cftsetpnumwidth{2.5em}
\cftsetrmarg{4em}
\usepackage{titlesec}
\titleformat{\section}{\Large\bfseries\raggedright}{\thesection}{0.7em}{}
\usepackage{fancyhdr}
\pagestyle{fancy}\fancyhf{}
\fancyhead[L]{\small\sffamily Open Problems Math | Research notebook}
\fancyhead[R]{\small 27 September 2026}
\fancyfoot[C]{\thepage}
\setlength{\parindent}{0pt}
\setlength{\parskip}{4pt plus 1pt}
\setlength{\emergencystretch}{3em}
\setcounter{tocdepth}{1}
\setcounter{secnumdepth}{1}
\setlist[itemize]{leftmargin=3em,itemsep=3pt,topsep=4pt,parsep=0pt}
\allowdisplaybreaks
\newcommand{\N}{\mathbb{N}}
\newcommand{\Z}{\mathbb{Z}}
\newcommand{\Q}{\mathbb{Q}}
\newcommand{\R}{\mathbb{R}}
\newcommand{\C}{\mathbb{C}}
\newcommand{\F}{\mathbb{F}}
\newcommand{\A}{\mathbb{A}}
\newcommand{\PP}{\mathbb P}
\newcommand{\E}{\mathbb{E}}
\newcommand{\eps}{\varepsilon}
\newcommand{\dd}{\,\mathrm d}
\newcommand{\sref}[2]{\hyperlink{source:#1:#2}{[S#2]}}
\newcommand{\eref}[2]{\hyperlink{extra:#1:#2}{[E#2]}}
% Turn the next switch off for a shorter reading copy. The quotations preserve
% every exactTarget field, including qualifications omitted from a short summary.
\newif\ifshowcatalogue
\showcataloguetrue
\newcommand{\cataloguescope}[1]{\ifshowcatalogue
 \paragraph{Exact scope in the supplied catalogue.}
 \begin{quote}\small #1\end{quote}\fi}
\usepackage{etoolbox}
\AtBeginEnvironment{itemize}{\small}
% Standalone entry; original section content is delimited for verification.
\begin{document}
\setcounter{section}{0}
%% BEGIN ORIGINAL SECTION CONTAINER
% ================================================================
% problemId: 3402
% Canonical problem ID: 3402
% exactTarget SHA-256: 537535a19e9ab22b3074cb8ff70bc305e7854d52ae180d24a321600b3ca4e027
\clearpage
\section*{\texorpdfstring{P vs. PSPACE}{P vs. PSPACE}}\label{3402}
\subsection{Definitions and mathematical statement}
For an input of length $n$, time counts computation steps and space counts work-tape cells. Define
\[
 \mathsf P=\bigcup_{k\ge1}\mathsf{DTIME}(n^k),\qquad
 \mathsf{PSPACE}=\bigcup_{k\ge1}\mathsf{DSPACE}(n^k).
\]
Machines must halt and decide their language correctly on every input. Determine whether these classes coincide. An algorithm with polynomial space and exponential time does not establish membership in $\mathsf P$.
\par\smallskip\textit{Formulation and concept references:} \sref{3402}{1}.
\cataloguescope{Determine whether P = PSPACE, that is, whether every language decidable by a deterministic Turing machine in polynomial space is decidable by a deterministic Turing machine in polynomial time.}
\subsection{Short English statement}
Can every decision problem solvable with polynomially much memory also be solved in polynomially much time?
%% BEGIN RESEARCH INSERTION
\subsection{Research attempt}
\paragraph{Current frontier.}
\textit{Research check: 27 September 2026.}
Williams's time-to-space simulation gives
$\mathsf{TIME}[t]\subseteq\mathsf{SPACE}[\sqrt{t\log t}]$ in the
multitape setting of Theorem 1.1, with the stated constructibility and
$t\geq n$ conventions. Its fixed-exponent consequences do not separate
the unions of all polynomial time and all polynomial space.
\sref{3402}{2}. The 2026 MFCS account discusses this development, not a
proof that the two classes in this entry differ. \sref{3402}{3}.
We investigate whether compact iteration of a transition map can remove
the exponential time of a polynomial-space computation, and then test
whether the new space simulation can be bootstrapped instead.

\paragraph{Attack route.}
A configuration graph has exponentially many vertices but a short local
transition rule. Repeated squaring is promising only if the iterated
transition maps stay in an efficiently constructible representation.
Alternatively, stronger uniform time-to-space compression could combine
with the space hierarchy. The first route needs a closure theorem under
composition; the second needs a substantially stronger simulation than
the theorem just quoted. We keep the two missing statements separate.

\paragraph{Attempt: an exact succinct-iteration target.}
Define a total decision problem $\mathcal I$. An instance consists of a
Boolean circuit $C$ computing $f_C:\{0,1\}^w\to\{0,1\}^w$, an initial
state $z$, a nonnegative integer $T$ in binary, and a coordinate $i$.
The answer is $(f_C^T(z))_i$; malformed instances are rejected. Circuit
encodings explicitly list inputs and outputs, so $w$ is bounded by the
input length. This problem belongs to $\mathsf{PSPACE}$: keep the current
$w$-bit state, a $\log(T+1)$-bit counter, and enough workspace to evaluate
$C$. The loop halts after $T$ iterations, even when $T$ is exponential
in its encoded length.

Conversely, let a fixed deterministic machine $M$ decide a language in
space $n^k$. Encode its configurations on input $x$ in $w=n^{O(1)}$
bits, using extra padding to include all state, head, and work-tape
information. A circuit $C_x$ of size polynomial in $w+|x|$ computes one
transition; invalid encodings can map to themselves. Make accepting and
rejecting configurations absorbing and add an acceptance flag coordinate.
The start configuration $z_x$ is explicitly computable. A halting
deterministic computation cannot revisit a nonhalting configuration, so
it halts within $2^w$ steps. Therefore
\[
 x\in L(M)\quad\Longleftrightarrow\quad
 (f_{C_x}^{\,2^w}(z_x))_{\rm accept}=1. \tag{14.1}
\]
The binary encoding of $2^w$ has only $w+1$ bits. This is a polynomial-time
reduction, with $M$ fixed and $x$ varying. It proves directly that
\[
 \mathcal I\in\mathsf P\quad\Longleftrightarrow\quad
 \mathsf P=\mathsf{PSPACE}. \tag{14.2}
\]
This standard configuration idea is used here as an explicit test bed,
not as an independently easier reformulation. The surrounding
configuration and simulation model is discussed in \sref{3402}{2},
Sections 1 and 5.

\paragraph{Attempt: a representation in which iteration does work.}
For the special case $f(z)=Az+b$ over $\F_2$, augment the state by $1$ and
put
\[
 \widetilde A=\begin{pmatrix}A&b\\0&1\end{pmatrix}.
 \qquad
 \begin{pmatrix}f^T(z)\\1\end{pmatrix}
 =\widetilde A^{\,T}\begin{pmatrix}z\\1\end{pmatrix}. \tag{14.3}
\]
Repeated squaring uses $O(\log(T+1))$ matrix multiplications, each with
$O((w+1)^3)$ field operations. All coefficients are bits. Thus affine
succinct iteration is polynomial-time in the \emph{encoded} parameters,
including binary $T$. This really does handle exponential iteration
counts, but only because the representation has a dimension closed
under composition.

For general circuits, the same notation hides duplication. If $C_j$
computes $f^{2^j}$, constructing $C_{j+1}$ by substituting $C_j$ into a
second copy gives a size recurrence of the form
$S_{j+1}\leq2S_j+O(w)$. This construction can grow as $2^j$, rather than
polynomially in $j$. Merely drawing two calls to the same subroutine is
not a Boolean circuit of size one call: they are evaluated on different
inputs. Fan-out shares a value, not a function evaluated at another
argument.

A degree bound also fails as a universal invariant. Introduce registers
$z,a_1,\ldots,a_r,t$, with a clock $t$ selecting the next update, and
perform $z\leftarrow za_t$. Starting at $z=1$, after $r$ updates one has
$z=a_1\cdots a_r$. Each selected update is quadratic in the data, but the
multilinear degree of its composition is $r$. This does not prove large
circuit size: the product has a small circuit. It does show why bounded
degree alone cannot justify an affine-style closure argument.

\paragraph{Attempt: the hierarchy route and its quantifiers.}
Here is a sufficient stronger simulation, written as a conditional
statement. Suppose that for every $\varepsilon>0$ and every suitable
time bound $t\geq n$,
\[
 \mathsf{TIME}[t]\subseteq\mathsf{SPACE}[t^{\varepsilon}]. \tag{14.4}
\]
Then $\mathsf P\ne\mathsf{PSPACE}$. Indeed, by the deterministic space
hierarchy choose
$L\in\mathsf{SPACE}[n^2]\setminus\mathsf{SPACE}[n]$.
If $\mathsf P=\mathsf{PSPACE}$, this same $L$ has time $O(n^k)$ for some
fixed $k$. Choose $0<\varepsilon<1/k$ in (14.4); constants and the
constructible time bound can be enlarged harmlessly. The resulting
space bound is $O(n^{k\varepsilon})\subseteq O(n)$, a contradiction.
This quantifier pattern is the strengthened-simulation direction
identified in \sref{3402}{2}, Section 5. It does not follow by applying
Theorem 1.1 repeatedly to its own output.

\paragraph{Stress test.}
To expose the failed bootstrap, let the first simulator use
$s=O(\sqrt{t\log t})$ work cells. A halting computation with this much
space need only have the crude configuration-count time bound
$2^{O(s+\log n)}$. Applying the time-to-space theorem to \emph{that}
implementation charges its new running time, not the original $t$.
Substituting the crude bound gives an exponentially large quantity, not
$t^{1/4}$. No polynomial-time preservation theorem for the first
simulator has been supplied. This is an accounting obstruction to the
proposed inference, not a lower bound on the best possible simulator.

Likewise, separation of $\mathsf{SPACE}[n]$ from a particular
$\mathsf{TIME}[n^a]$ does not separate $\mathsf{PSPACE}$ from every
polynomial time bound. The exponent of a putative decider is unrestricted.
The affine algorithm cannot process the nonlinearity of an arbitrary
configuration circuit, and finite experiments on randomly chosen maps
do not address the worst-case uniform circuit family in (14.1).
A conditional argument using $\mathsf P\ne\mathsf{NP}$ would already
settle this entry negatively because
$\mathsf P\subseteq\mathsf{NP}\subseteq\mathsf{PSPACE}$; it would not
be an unconditional separation.

\paragraph{Outcome.}
\textbf{Outcome: Exploratory approach with unresolved gap.}
The notebook supplies an exact succinct-iteration reduction, a genuine
polynomial-time affine special case, and a quantified stronger-simulation
condition that would separate the classes. Neither efficient closure of
general transition circuits nor (14.4) is established. The attempted
bootstrap of square-root-space simulation fails because it reuses the
wrong time parameter. These auxiliary deductions are not claimed to be
historically new or to resolve the original class equality.
%% END RESEARCH INSERTION
\subsection{Sources}
\begin{itemize}\raggedright
\item[\textnormal{[S1]}]\hypertarget{source:3402:1}{} Walter S. Brainerd and Lawrence H. Landweber et al., Computability, Decidability, Complexity, third edition.
\url{https://art.torvergata.it/bitstream/2108/32769/60437/ComputabilityDecidabilityComplexity.pdf}.
{\small\textit{Catalogue formulation source.}}
\item[\textnormal{[S2]}]\hypertarget{source:3402:2}{} Simulating Time With Square-Root Space — Ryan Williams, February 24, 2025.
\url{https://eccc.weizmann.ac.il/report/2025/017/download}.
{\small\textit{Catalogue research/status reference; not a proof endorsement.}}
\item[\textnormal{[S3]}]\hypertarget{source:3402:3}{} Some Recent Developments in Space Complexity — R. Ryan Williams, MFCS 2026.
\url{https://drops.dagstuhl.de/entities/document/10.4230/LIPIcs.MFCS.2026.3}.
{\small\textit{Catalogue research/status reference; not a proof endorsement.}}
\item[\textnormal{[S4]}]\hypertarget{source:3402:4}{} Some conditions implying if P=NP then P=PSPACE — Ismael Rodriguez (February 10, 2026).
\url{https://arxiv.org/abs/2602.10073}.
{\small\textit{Catalogue research/status reference; not a proof endorsement.}}
\end{itemize}

%% END ORIGINAL SECTION CONTAINER
\end{document}
